\documentclass[11pt]{article}
\usepackage[margin=1in]{geometry}
\usepackage{times}
\usepackage{hyperref}
\hypersetup{colorlinks=true, linkcolor=black, urlcolor=blue, citecolor=black}
\title{Symmetry Does Not Care Who Denies It: Stefan Marinov's Asymmetric Homopolar Claims}
\author{Flyxion \\ \small Independent Researcher}
\date{}
\begin{document}
\maketitle

\begin{abstract}
Stefan Marinov, a Bulgarian physicist, promoted from the 1970s until his death in 1997 a series of devices, most notably the "Siberian Coliu" and various asymmetric homopolar motor-generator configurations, which he claimed demonstrated violations of both energy conservation and the relativity of motion, framing these results as evidence against Einstein's special relativity more broadly. This paper argues that Marinov's specific electromagnetic asymmetry claims trace to well-documented and independently reproducible measurement artifacts in homopolar device wiring and brush contact configurations, that his broader anti-relativity program has been independently and specifically rebutted on its own theoretical terms by physicists who engaged directly with his published derivations, and that Marinov's death, by suicide following years of what he described as academic ostracism, does not itself constitute evidence for or against the correctness of claims that were addressed and found wanting through ordinary technical review well before his death.
\end{abstract}

\section{Introduction}

Stefan Marinov held a physics doctorate and worked briefly in academic positions before being dismissed amid disputes over his increasingly unconventional physics claims, subsequently self-publishing extensively and building demonstration devices, most prominently variations on the homopolar generator, a disc or asymmetric conductor rotating in a magnetic field, which he claimed showed measurable asymmetries in induced current or output voltage depending on which reference frame or measurement configuration was used, asymmetries he interpreted as evidence against special relativity's postulate that physical laws take the same form in all inertial frames, and in a related but distinct claim, as evidence of net energy production from certain rotating configurations.

\section{The Measurement Artifact in Homopolar Asymmetry Claims}

Homopolar generators and motors have a long documented history of measurement subtleties related to exactly where sliding brush contacts are placed relative to the rotating conductor and the surrounding stationary frame, since the induced electromotive force in a homopolar device depends on the specific circuit path completed by the external measurement leads in a way that is well understood within conventional electromagnetism but has repeatedly produced apparent paradoxes for experimenters who do not fully account for it, a category of well-known conceptual and measurement pitfall predating Marinov's work by decades and discussed extensively in the electromagnetism pedagogical literature independent of any relativity-related claim. Independent physicists who examined Marinov's specific device configurations, including detailed published critiques in mainstream physics venues during his lifetime, identified specific brush placement and circuit topology choices in his apparatus consistent with this well-known artifact category as the source of his reported asymmetries, rather than any genuine departure from either conventional electromagnetism or special relativity.

\section{Why the Anti-Relativity Program Was Independently and Substantively Engaged**}

Unlike several claims surveyed elsewhere in this series that circulate largely without direct engagement from the mainstream physics community, Marinov's specific theoretical arguments against special relativity received direct, substantive published rebuttal from physicists who engaged with his actual derivations on their own mathematical terms, identifying specific errors in how he applied Lorentz transformations and defined simultaneity in his proposed experiments, rather than dismissing the claims without engagement; this matters because it means the relevant scientific response was not merely disinterest or institutional gatekeeping but active, specific technical correction that Marinov's published responses did not successfully address on the terms his critics had raised.

\section{The Biographical Framing and Why It Does Not Bear on the Physics}

Marinov died in 1997, by suicide, after describing in his own writings a long period of professional isolation and frustration at what he experienced as the physics community's refusal to seriously engage with his work, a genuinely difficult personal history that has since been cited by some supporters as evidence of unjust suppression of correct physics. The specific technical rebuttals of Marinov's homopolar measurement claims and relativity derivations, discussed above, were published and addressed his actual mathematics and apparatus configuration directly, a form of engagement inconsistent with simple institutional dismissal, and the tragic circumstances of his death, while worth acknowledging as a serious human cost, do not alter the independent technical content or validity of either his claims or the rebuttals to them, which must be evaluated on their own terms regardless of the emotional weight the biographical narrative carries.

\section{The Entropic Gravity Framing}

Marinov's claims concern electromagnetic asymmetry and special relativity rather than gravitation directly, and are included in this series as a comparative case in the broader "conservation law violation" and "mainstream physics is fundamentally wrong" genre; the entropic account of gravitation addressed throughout this series is explicitly built to reduce to standard general relativity in appropriate limits rather than to overturn special or general relativity's confirmed predictions, placing it in a different category from Marinov's program, which specifically targeted and failed to substantiate a departure from those same confirmed predictions.

\section{Objections}

Supporters argue that mainstream physics has a structural bias against claims threatening foundational theories regardless of their merit, making the negative reception of Marinov's work unsurprising independent of whether his physics was correct. Structural resistance to genuinely revolutionary claims is a real phenomenon worth taking seriously in the abstract, but it does not substitute for engagement with the specific technical rebuttals Marinov's homopolar and relativity claims received, which addressed his actual apparatus and mathematics directly rather than dismissing the work on reputational grounds alone, and which his own published responses did not resolve.

\section{Conclusion}

Marinov's asymmetric homopolar claims trace to a well-documented category of measurement artifact predating his own work, and his broader anti-relativity program received direct, substantive technical rebuttal engaging his specific derivations rather than institutional dismissal, a combination that leaves his central claims unsupported on technical grounds independent of, and prior to, the genuinely difficult personal and professional circumstances of his later life.

\end{document}
